% circuit-sim-lab-report — report.tex template
% Course 0714-EEE-2102, Circuit Simulation Laboratory
% Copy to <expN>/Lab_report_latex/report.tex, set the experiment number
% below, fill the header block, and adapt the sections.
%
% Compile (twice, from Lab_report_latex/):
%   env pdflatex -interaction=nonstopmode report.tex

\documentclass[12pt,a4paper]{article}

% ----- Page geometry: A4, 1 inch margins all sides -----
\usepackage[margin=1in]{geometry}

% ----- Times New Roman (text + math) -----
\usepackage{newtxtext,newtxmath}

% ----- Spacing: 1.5 line spacing, 6 pt paragraph spacing -----
\usepackage{setspace}
\onehalfspacing
\usepackage{parskip}          % paragraph gaps instead of indents
\setlength{\parskip}{6pt plus 1pt}

% ----- Math, SI units -----
\usepackage{amsmath}
\usepackage{siunitx}

% ----- Graphics, tables, captions, floats -----
\usepackage{graphicx}
\usepackage{booktabs}
\usepackage{caption}
\usepackage{float}
% Figures live in <expN>/Graphs/; add ./figures/ for report-local images
% (e.g. schematic exports).
\graphicspath{{../Graphs/}{./figures/}}

% ----- Lists: compact, guide-style itemization -----
\usepackage{enumitem}
\setlist{nosep, leftmargin=2em, label=\textbullet}

% ----- Front matter and PDF output -----
\usepackage{pdfpages}
\usepackage[hidelinks]{hyperref}

% ----- Headings: bold, left-aligned, 14/15 pt -----
\usepackage{titlesec}
\titleformat{\section}
  {\normalfont\Large\bfseries}{\thesection}{0.8em}{}
\titleformat{\subsection}
  {\normalfont\large\bfseries}{\thesubsection}{0.8em}{}
% \Large = 14.4 pt, \large = 14.4→~14 pt at 12pt base; matches 14–16 pt rule
\titlespacing*{\section}{0pt}{*3}{*2}
\titlespacing*{\subsection}{0pt}{*2.5}{*1.5}

% ----- Captions: centered, 11 pt, figures below / tables above -----
\captionsetup{font=small, labelfont=bf, justification=centering}
\captionsetup[table]{position=above, skip=4pt}
\captionsetup[figure]{position=below, skip=4pt}

% ----- Page number: bottom-center -----
\usepackage{fancyhdr}
\pagestyle{fancy}
\fancyhf{}
\fancyfoot[C]{\thepage}
\renewcommand{\headrulewidth}{0pt}

% ----- Experiment number: sections then number N.1–N.10 -----
% Guideline numbering: "1.1 Objective" ... "1.10 Reference" where the
% prefix is the experiment number. Figures, tables, and equations keep
% plain consecutive numbers (Figure 1, Table 1, (1)) per the guideline.
\newcommand{\expnumber}{1}   % <-- CHANGE to the assigned experiment number
\renewcommand{\thesection}{\expnumber.\arabic{section}}

\begin{document}

% =====================================================================
% FRONT MATTER (pages 1–2), per the instructor's Intructions.txt:
%   page 1: Cover_Page.pdf    page 2: 2nd_page.pdf
% \includepdf adds a separate page and resets the page counter by
% default; keep the body starting at page 1 of its own numbering.
% =====================================================================
\includepdf[pages={1}]{../Cover_Page.pdf}
\includepdf[pages={1}]{../2nd_page.pdf}

% =====================================================================
% HEADER BLOCK
% =====================================================================
\begin{center}
  {\large\bfseries Experiment No. \expnumber}\\[4pt]
  {\large\bfseries Name of the Experiment: Simulation, Analysis, and
    Evaluation of Transient and Steady-State Responses of RC Circuits
    with Switching, Pulse, and Sinusoidal Inputs Using PSpice}
\end{center}
\vspace{6pt}

% =====================================================================
\section{Objective}
% 2–4 items, action verbs (design/simulate/analyze/verify/investigate),
% name the circuit/software concepts. No theory here.
% =====================================================================
\begin{enumerate}[label=\roman*.]
  \item To design and simulate RC circuits with switching, pulse, and
        sinusoidal inputs using PSpice.
  \item To analyze the transient and steady-state responses of the
        simulated circuits.
  \item To compare simulated and theoretical time constants through
        percentage-error analysis.
\end{enumerate}

% =====================================================================
\section{Theory}
% Principle of operation, governing equations (numbered), parameter
% definitions, expected waveform behavior, assumptions. Explain why
% the circuit behaves as observed — not a textbook dump.
% =====================================================================
An RC circuit consists of a resistor $R$ in series with a capacitor
$C$... [principle of operation]

The charging of the capacitor through $R$ is governed by
\begin{equation}
  v_C(t) = V_s \left(1 - e^{-t/\tau}\right), \qquad \tau = RC
  \label{eq:charging}
\end{equation}
where $V_s$ is the source voltage and $\tau$ the time constant...

% =====================================================================
\section{Design}
% Include ONLY if the experiment requires component/parameter design
% or selection. Otherwise delete this section entirely.
% =====================================================================
... [circuit specifications, design equations, component-value
calculations, selected values, constraints]

% =====================================================================
\section{List of Components}
% Table of all components, instruments, and software used.
% =====================================================================
\begin{table}[H]
  \centering
  \caption{List of components, instruments, and software}
  \label{tab:components}
  \begin{tabular}{clcc}
    \toprule
    Sl. No. & Component/Instrument & Specification/Value & Quantity \\
    \midrule
    1 & Resistor           & \SI{1}{\kilo\ohm}     & 1 \\
    2 & Capacitor          & \SI{100}{\micro\farad} & 1 \\
    3 & Voltage Source     & DC, \SI{10}{\volt}    & 1 \\
    4 & Switch             & Voltage-controlled    & 1 \\
    5 & Simulation Software & PSpice              & 1 \\
    \bottomrule
  \end{tabular}
\end{table}

% =====================================================================
\section{Working Procedure}
% Numbered steps, reproducible by another student.
% NO results, NO conclusions here.
% =====================================================================
\begin{enumerate}
  \item ...
\end{enumerate}

% =====================================================================
\section{Working Diagram}
% The ACTUAL simulation schematic export/screenshot — not a redrawn
% drawing. Save it under Graphs/ or ./figures/ and reference by bare
% filename (graphicspath covers both).
% =====================================================================
\begin{figure}[H]
  \centering
  \includegraphics[width=0.8\textwidth]{schematic.png}
  \caption{PSpice schematic of the RC circuit under simulation}
  \label{fig:schematic}
\end{figure}

% =====================================================================
\section{Calculation}
% Given values → equations → substitution → intermediate steps →
% final value with units. Show the steps, not just the answer.
% =====================================================================
Theoretical time constant with $R = \SI{1}{\kilo\ohm}$ and
$C = \SI{100}{\micro\farad}$:
\begin{equation}
  \tau = RC
       = \left(\SI{1}{\kilo\ohm}\right)\left(\SI{100}{\micro\farad}\right)
       = \SI{0.1}{\second}
  \label{eq:tau-calc}
\end{equation}

% =====================================================================
\section{Result Analysis}
% Waveforms, simulated vs. theoretical comparison tables, % error,
% parameter-variation results. Every figure/table numbered, captioned,
% and referred to in the prose.
% =====================================================================
\begin{figure}[H]
  \centering
  \includegraphics[width=0.85\textwidth]{exp1_1.png}
  \caption{Transient charging response of the RC circuit}
  \label{fig:charging}
\end{figure}

From Figure~\ref{fig:charging}, the simulated time constant is ...

\begin{table}[H]
  \centering
  \caption{Simulated vs. theoretical time constants}
  \label{tab:tau-compare}
  \begin{tabular}{lccc}
    \toprule
    Quantity & Theoretical & Simulated & \% Error \\
    \midrule
    $\tau$ (charging) & \SI{0.1}{\second} & \SI{0.098}{\second} & 2.0 \\
    $\tau$ (discharging) & \SI{0.1}{\second} & \SI{0.101}{\second} & 1.0 \\
    \bottomrule
  \end{tabular}
\end{table}

% =====================================================================
\section{Discussion}
% Analysis, NOT a repeat of Result Analysis: why waveforms have this
% shape, why simulated differs from theoretical, parameter effects,
% model limitations, practical significance.
% =====================================================================
...

% =====================================================================
\section{Reference}
% IEEE style, numbered in order of first citation in the text.
% =====================================================================
\begin{enumerate}[label={[\arabic*]}, leftmargin=2.5em]
  \item W. H. Hayt, J. E. Kemmerly, and S. M. Durbin, \textit{Engineering
        Circuit Analysis}, 8th ed. New York, NY, USA: McGraw-Hill, 2012.
\end{enumerate}

\end{document}
